% from ega3-1-fr.tex
\oldpage[0\textsubscript{III}]{17}
Pour voir que la condition est suffisante, appliquons le principe de
récurrence noethérienne à l'ensemble $\mathfrak{M}$ des parties fermées $Y$
de $X$ telles que la restriction $h|Y$ soit constructible~; on peut donc
supposer que pour toute partie fermée $Y\ne X$ de $X$, $h|Y$ est
constructible. Si $X$ n'est pas irréductible, la restriction de $h$ à
chacune des composantes irréductibles $X_i$ de $X$ (en nombre fini) est donc
constructible, et il résulte alors aussitôt de la définition
(\hyperref[0.9.3.1-fr]{9.3.1}) que $h$ est constructible. Si $X$ est
irréductible, il existe par hypothèse une partie ouverte non vide $U$ de $X$
dans laquelle $h$ est constante~; d'autre part, la restriction de $h$ à
$X\setminus U$ est constructible par hypothèse, et il en résulte aussitôt
que $h$ est constructible.

\begin{corollary}[9.3.3]
\label{0.9.3.3-fr}
Soit $X$ un espace noethérien dans lequel toute partie fermée irréductible
admet un point générique. Si $h$ est une application de $X$ dans un ensemble
$T$ telle que, pour tout $t\in T$, $h^{-1}(t)$ soit constructible, alors $h$
est constructible.
\end{corollary}

En effet, si $Y$ est une partie fermée irréductible de $X$ et $y$ son point
générique, $Y\cap h^{-1}(h(y))$ est constructible et contient $y$, donc
(\hyperref[0.9.2.2-fr]{9.2.2}) cet ensemble contient une partie ouverte non
vide de $Y$, et il suffit d'appliquer (\hyperref[0.9.3.2-fr]{9.3.2}).

\begin{proposition}[9.3.4]
\label{0.9.3.4-fr}
Soient $X$ un espace noethérien dans lequel toute partie fermée irréductible
admet un point générique, $h$ une application constructible de $X$ dans un
ensemble ordonné. Pour que $h$ soit semi-continue supérieurement dans $X$,
il faut et il suffit que pour tout $x\in X$ et toute générisation
(\hyperref[0.2.1.2-fr]{$0_{\mathrm I}$, 2.1.2}) $x'$ de $x$, on ait
$h(x')\leq h(x)$.
\end{proposition}

La fonction $h$ ne prend qu'un nombre fini de valeurs~; dire qu'elle est
semi-continue supérieurement signifie donc que pour tout $x\in X$,
l'ensemble $E$ des $y\in X$ tels que $h(y)\leq h(x)$ est un voisinage de $x$.
Par hypothèse, $E$ est une partie constructible de $X$~; d'autre part, dire
qu'une partie fermée irréductible $Y$ de $X$ contient $x$ signifie que son
point générique $y$ est une générisation de $x$~; la conclusion résulte
alors de (\hyperref[0.9.2.5-fr]{9.2.5}).

\section{Compléments sur les modules plats}
\label{section:0.10-fr}

Pour les démonstrations des propriétés énoncées sans démonstration dans les
n\textsuperscript{os} (\hyperref[subsection:0.10.1-fr]{10.1}) et
(\hyperref[subsection:0.10.2-fr]{10.2}), nous renvoyons le lecteur à Bourbaki,
\emph{Alg. comm.}, chap. II et III.

\subsection{Relations entre modules plats et modules libres}
\label{subsection:0.10.1-fr}

\begin{env}[10.1.1]
\label{0.10.1.1-fr}
Soient $A$ un anneau, $\mathfrak{J}$ un idéal de $A$, $M$ un $A$-module~; pour
tout entier $p\geq0$, on a un homomorphisme canonique de
$(A/\mathfrak{J})$-modules
\[
\label{0.10.1.1.1-fr}
  \varphi_p:(M/\mathfrak{J}M)\otimes_{A/\mathfrak{J}}
  (\mathfrak{J}^p/\mathfrak{J}^{p+1})
  \longrightarrow \mathfrak{J}^pM/\mathfrak{J}^{p+1}M
  \tag{10.1.1.1}
\]
qui est évidemment \emph{surjectif}. Nous désignerons par
$\operatorname{gr}(A)=\bigoplus_{p\geq0}\mathfrak{J}^p/\mathfrak{J}^{p+1}$ l'anneau gradué
associé à $A$ filtré par les $\mathfrak{J}^p$, par
$\operatorname{gr}(M)=\bigoplus_{p\geq0}\mathfrak{J}^pM/\mathfrak{J}^{p+1}M$ le
$\operatorname{gr}(A)$-module gradué associé à $M$ filtré par les $\mathfrak{J}^pM$~; on a
donc $\operatorname{gr}_p(A)=\mathfrak{J}^p/\mathfrak{J}^{p+1}$,
$\operatorname{gr}_p(M)=\mathfrak{J}^pM/\mathfrak{J}^{p+1}M$~; les $\varphi_p$ définissent
un homomorphisme \emph{surjectif} de $\operatorname{gr}(A)$-modules gradués
\[
\label{0.10.1.1.2-fr}
  \varphi:\operatorname{gr}_0(M)\otimes_{\operatorname{gr}_0(A)}
  \operatorname{gr}(A)\longrightarrow\operatorname{gr}(M).
  \tag{10.1.1.2}
\]
\end{env}

